% Slide 11 -- The geometric trick: parameterizing the circle
\begin{frame}{The geometric trick: parameterizing the articular circle}
  \begin{columns}[T]
    \begin{column}{0.55\textwidth}
      \begin{block}{From a circle to a 1D angular map}
        Each per-slice circle has an in-plane basis $(\mathbf{u},\mathbf{v})$ and center $\mathbf{c}$.
        Any inlier $\mathbf{p}$ has angular coordinate
        \[
          \theta = \mathrm{atan2}\big((\mathbf{p}-\mathbf{c})\!\cdot\!\mathbf{v},\; (\mathbf{p}-\mathbf{c})\!\cdot\!\mathbf{u}\big).
        \]
      \end{block}

      \begin{block}{Why this is powerful}
        The articular circle becomes a \alert{1D tendon map}:
        each $\theta$ on the circle corresponds to a known anatomical direction at the humeral head.
      \end{block}
    \end{column}
    \begin{column}{0.45\textwidth}
      \centering
      \begin{tikzpicture}[scale=1.0]
        \filldraw[medSoft] (0,0) circle (2.0);
        \draw[thick, medBlue] (0,0) circle (2.0);
        \draw[->] (-2.4,0) -- (2.4,0) node[right, font=\tiny] {$\mathbf{u}$};
        \draw[->] (0,-2.4) -- (0,2.4) node[above, font=\tiny] {$\mathbf{v}$};
        \filldraw[medBlue] (0,0) circle (1.5pt);
        \node[medBlue, font=\tiny, below left] at (0,0) {$\mathbf{c}$};
        \filldraw[medAccent] (60:2.0) circle (1.8pt);
        \draw[medAccent, dashed] (0,0) -- (60:2.0);
        \node[medAccent, font=\tiny, above right] at (60:2.0) {$\mathbf{p}$};
        \draw[medAccent] (0.6,0) arc (0:60:0.6);
        \node[medAccent, font=\tiny] at (35:0.85) {$\theta$};
      \end{tikzpicture}
    \end{column}
  \end{columns}
\end{frame}

% Slide 11b -- Anatomical sector mapping
\begin{frame}{Angular sectors $\to$ rotator-cuff tendons}
  \begin{columns}[T]
    \begin{column}{0.5\textwidth}
      \begin{block}{Anatomical sector mapping}
        Tendon insertions fall in \alert{predictable sectors}:
        \begin{itemize}
          \item \textbf{Subscapularis}: $\sim 0$--$45^\circ$.
          \item \textbf{Supraspinatus}: $\sim 45$--$135^\circ$.
          \item \textbf{Infraspinatus}: $\sim 135$--$225^\circ$.
          \item \textbf{Teres minor}: $\sim 225$--$270^\circ$.
        \end{itemize}
      \end{block}

      \begin{block}{Localization rule}
        A focal drop in inlier ratio in a sector $\to$ \alert{candidate tear} in the corresponding tendon, at the slice where the drop occurs.
      \end{block}
    \end{column}
    \begin{column}{0.5\textwidth}
      \centering
      \begin{tikzpicture}[scale=0.85]
        \filldraw[medSoft] (0,0) circle (2.2);
        \draw[thick, medBlue] (0,0) circle (2.2);
        \filldraw[medTeal!30] (0,0) -- (2.2,0) arc (0:45:2.2) -- cycle;
        \filldraw[medTeal!50] (0,0) -- (45:2.2) arc (45:135:2.2) -- cycle;
        \filldraw[medTeal!30] (0,0) -- (135:2.2) arc (135:225:2.2) -- cycle;
        \filldraw[medTeal!50] (0,0) -- (225:2.2) arc (225:270:2.2) -- cycle;
        \node[font=\tiny] at (22:1.5) {Subscap.};
        \node[font=\tiny] at (90:1.5) {Supraspin.};
        \node[font=\tiny] at (180:1.5) {Infraspin.};
        \node[font=\tiny] at (247:1.5) {T. minor};
        \draw[medAccent, ultra thick] (60:2.0) -- (60:2.6);
        \draw[medAccent, ultra thick] (75:2.6) -- (75:2.0);
        \node[medAccent, font=\tiny, above] at (67:2.7) {focal drop};
      \end{tikzpicture}
      \\[0.3em]
      {\footnotesize The articular circle becomes a 1D \emph{tendon map}.}
    \end{column}
  \end{columns}
\end{frame}
